\documentclass[../main.tex]{subfiles}
\graphicspath{{\subfix{../img/}}}
\begin{document}
	Un sistema intelligente deve essere in grado di prendere decisioni in
	modo autonomo, in funzione delle condizioni in cui si trova ad
	operare.
	\subsection{Test di Turing}
	Il test di Turing nasce nel 1950 come test per capire se una macchina si possa veramente considerare intelligente.\\ \\
	\textbf{Condizione}: "Un computer è da considerarsi intelligente se, nell'interazione a distanza con un essere umano, la persona non è in grado di distinguere se sta interagendo con un uomo o una macchina".\\ \\
	Test di Turing più "importante" è stato fatto nel 1991, dove dei giudici facendo delle domande dovevano capire se dall'altra parte ci fosse una persona o una macchina.\\
	Nonostante la tecnologia non fosse ancora avanzata e la macchina rispondesse con un contesto sbagliato ad alcune domande, diversi giudici sono comunque stati ingannati.\\ \\
	Questo portò alla nascita di un dibattito sulla possibilità che le macchine potessero avere una coscienza o che, semplicemente, imitassero a "pappagallo" le indicazioni date dall'umano al momento della programmazione.
	\subsection{La stanza cinese}
	Nasce nel 1980\\
	\begin{itemize}
		\item Supponendo che il sistema superi il test di Turing, il computer prende simboli cinesi in entrata, esegue un programma e produce altri simboli.
		\item Nella stanza cinese c’è una persona a cui viene dato un libro con la versione inglese del programma utilizzato.
		\item La persona può elaborare i testi in cinese in arrivo, seguendo le istruzioni, ma non capisce realmente i simboli.
		\item Allo stesso modo, nemmeno la macchina riesce a comprendere ciò che legge.
	\end{itemize}
	Il computer è un manipolatore di simboli, quindi non capisce quello che viene espresso in cinese tanto quanto una persona che parla solo inglese.
	\subsection{Al forte vs Al debole (o narrow vs general)}
	\begin{center}
		\begin{tabular}{ |p{8cm}|p{8cm}| }
			\hline
			\textbf{Al debole} & \textbf{Al forte} \\
			\hline
			Assumono le sembianze di un comportamento intelligente che può essere modellato e usato dai computer per risolvere problemi complessi & Ipotetica macchina con abilità umane cognitive \\
			\hline
			Si riferisce ad un sistema programmato per risolvere un ampio assortimento di problemi, ma opera con impostazioni predefinite & Con propria abilità di pensiero e che risolve task da sé \\
			\hline
			Non hanno mente propria & Macchine che hanno forte abilità umana cognitiva \\
			\hline
			Alexa e Siri sono i migliori esempi di Al debole & Un concetto ipotetico che non esiste ancora nella vita reale\\
			\hline
		\end{tabular}
	\end{center}
	\subsection{Oridine del termine}
	Espressione "Artificial Intelligence" coniata nel 1956 dal matematico John McCarthy-Logic durante un seminario.
	\subsubsection{Gli inizi (1952-1959)}
	\begin{figure}[h]
		\centering
		\adjustbox{cfbox=white 5pt 0cm}{\includegraphics[scale=0.3]{example-image-golden}}
		\caption{An example of a problem solved with. Evans' ANALOGY program (1968)}
		\label{fig:im-1}
	\end{figure}
	\begin{itemize}
		\item Per fare analogie faceva uso di una buona/adatta" rappresentazione dell'informazione
		\item II modo in cui qualcosa viene rappresentato su un particolare sistema di elaborazione. delle informazioni influisce in modo cruciale sul tipo di operazioni che possiamo eseguire in modo naturale e rapido
		\item Ciò che ha fatto capire l'importanza della rappresentazione dell'informazione e stata. la difficoltà tecnica di programmare i computer
	\end{itemize}
	\textbf{Programma Analogico di Evan}: Trasforma un'immagine in un grafo di conoscenza.
	\subsubsection{1968}
	\begin{itemize}
		\item Capacita di calcolo limitata 
		\item Restrizioni sul dominio (oggetti semplici rispetto al reale e poche relazioni).
		\item Programma:
		\begin{itemize}
			\item Tradurre il disegno di una rappresentazione vettoriale di forme geometriche 
			\item Evidenziare le relazioni fra le forme dei fatti noti (la sequenza delle immagini date) 
			\item Selezionare un insieme di immagini possibili secondo i vincoli del dominio 
			\item Scegliere una di queste immagini
		\end{itemize}
	\end{itemize}
	\subsubsection{1969 - 1979}
	\textbf{Sistemi esperti}: Gli informatici cercano di trasferire la conoscenza di un esperto, in un settore limitato, in un insieme di regole: if condition then action
	\subsubsection{1995 - 2000: Nascita del Web e delle ontologie}
	II grande successo del web ha portato in primo piano, ma in una forma nuova, la questione della rappresentazione simbolica della conoscenza.\\ \\
	Questioni più difficili sono l'interoperabilità delle applicazioni e la gestione delle risorse sulla rete.\\ \\
	II concetto fondamentale e quello di ontologia, che ha cominciato a precisarsi negli anni $'80$ Le tecnologie semantiche saranno essenziali per fornire una piattaforma tecnologica adeguata.
	\subsubsection{2011 - present}
	L'accesso a grandi quantità di dati ("big data"), computer più economici e veloci e tecniche avanzate di apprendimento automatico sono stati applicati con successo a molti problemi.\\
	I progressi nel deep learning e nella ricerca hanno portato a progressi nell'elaborazione di immagini e video, nell'analisi del testo e persino nel riconoscimento vocale.
	\subsection{Il gioco In IA}
	\begin{itemize}
		\item \textbf{Chess games}: Lo svolgersi del gioco può essere interpretato come un albero in cui la radice è la posizione di partenza e le foglie le posizioni finali.
		\item\textbf{ Gioco del GO}.
		\item \textbf{Jeopardy!}: è un quiz televisivo statunitense che consiste in una cultura generale della competizione tra i vari concorrenti, trasmesso dall'emittente NBC dal 1964.
		\item \textbf{Pluribus}: Gioco simile al Poker, è riuscito a replicare i comportamenti dell'essere umano, cioè il bluff;
		\item \textbf{Stratego}: Come racconta uno studio pubblicato su Science il 1 dicembre, l'AI - DeepNash ha imparato a bluffare con i pezzi meno di valore e a sacrificare quelli più pregiati pur di ottenere la vittoria e guadagnare informazioni sulla strategia dell'avversario.\\
		Queste scelte lo rendono imprevedibile, raggiungendo l'84\% di tasso di vincita durante 50 partite svolte online.
	\end{itemize}
	\begin{figure}[h]
		\centering
		\adjustbox{cfbox=white 5pt 0cm}{\includegraphics[scale=0.3]{example-image-golden}}
		\caption{Immagine di un Chess Game}
		\label{fig:im-2}
	\end{figure}
	\subsection{Ragionamento Deduttivo vs Ragionamento Induttivo}
	\textbf{ASI}: Intelligenza artificiale speculativa. ASI è dove le macchine sono autocoscienti, oltre le capacità e conoscenze umane.\\ \\
	\textbf{Approccio psicologico}: Obiettivo di creare una macchina pensante (progetto, per ora, fallimentare). \\ \\
	\textbf{Approccio ingegneristico}: Obiettivo di programmare una macchina che risolva problemi in modo più efficiente. \\ \\
	Per dare una coscienza alla macchina ci vorrebbe una tecnologia completamente diversa, non sarebbe possibile utilizzare quella computerizzata che utilizziamo adesso, serve un approccio completamente diverso.
	\subsection{Conoscenza e comportamento razionale}
	La conoscenza è informazione disponibile per un'azione razionale.\\ \\
	\textbf{Comportamento razionale}: Fare la cosa più giusta, quella da cui ci si aspetta il massimo risultato, a fronte delle informazioni disponibili.\\ \\
	Una corretta inferenza non sempre esaurisce il comportamento razionale (es. in certe situazioni va decisa anche in assenza di inferenze). Non tutti i comportamenti razionali implicano inferenze (es. togliere una mano dal fuoco).
	\subsection{Approccio simbolico vs approccio subsimbolico}
	Gli individui agiscono sulla conoscenza e sulle percezioni del mondo attraverso una attività di ragionamento:
	\begin{itemize}
		\item \textbf{Deduttivo}: Dalle premesse alle conclusioni
		\item \textbf{Abduttivo}: Dagli effetti alle possibili cause
		\item \textbf{Induttivo}: Da fatti specifici a regole generali
	\end{itemize}
	Deduttivo e abduttivo hanno un approccio simbolico, induttivo uno sub-simbolico.\\
	Questi ragionamenti servono per produrre una nuova conoscenza e, eventualmente, decidere una azione.\\ \\
	\textbf{Top Down o simbolico}: La conoscenza è rappresentata da simboli. Vengono utilizzate la logica e la capacità di fare inferenze.\\
	\textbf{Bottom UP o subsimbolico} (basato sul segnale): Non esiste una rappresentazione esplicita della conoscenza. Vengono utilizzate reti neurali, approcci evolutivi, ecc.\\ \\
	Non esiste una metodologia migliore dell’altra, dipende solo da che contesti vengono utilizzate.
	\subsubsection{Approccio simbolico}
	L’approccio simbolico alla costruzione di un AI si serve di una descrizione esplicita della conoscenza di un dominio e fa uso di un linguaggio formale (adatto ad essere utilizzato da un automa) per attingere a questa base di conoscenza e inferire nuova conoscenza (o azioni da intraprendere).\\
	Per poter realizzarli si sono sviluppati framework basati sulla logica come classe generale di rappresentazioni e linguaggio formale.
	\paragraph{Logica}
	La logica è la scienza che fornisce agli uomini gli strumenti indispensabili per verificare con sicurezza la correttezza del ragionamento.\\
	La logica fornisce gli strumenti formali per:
	\begin{itemize}
		\item \textbf{Esprimere} inferenze in termini di operazioni su espressioni simboliche
		\item \textbf{Dedurre} le conseguenze da determinate premesse
		\item \textbf{Studiare} la verità o la falsità su certe proposizioni data la verità o la falsità di altre proposizioni
		\item \textbf{Stabilire} la coerenza e la validità di una data teoria
	\end{itemize}
	\textbf{Ragionamento deduttivo}: Modo per inferire nuove conoscenze nella logica, se le premesse sono vere, allora le conclusioni sono vere.\\
	\textbf{Regole di inferenza}: Composte da due parti divise dal simbolo “=>” che si legge come “è possibile derivare”. \\ \\
	Il \textbf{sillogismo} è un tipico esempio di ragionamento deduttivo. \\
	In generale, in logica si studia la relazione tra le formule di un sistema logico-simbolico in cui:
	\begin{itemize}
		\item Il linguaggio delle formule è definito con precisione
		\item Il significato delle formule è stabilito in modo univoco
	\end{itemize}
	E’ possibile trovare un'affermazione falsa partendo da due proposizioni vere utilizzando un linguaggio naturale (che per sua natura è ambiguo).\\
	\textbf{Logica moderna}: Rappresenta in modo formale la corrispondenza tra espressioni linguistiche e significato inteso.\\
	\textbf{Formalizzazione (astrazione)}: Comprensione delle informazioni, con ampia perdita di dettagli (sfumature) e guadagno di precisione (non tutte le espressioni linguistiche sono “buoni candidati” per questa traduzione).\\ \\
	Uno dei compiti più famosi nello studio del ragionamento deduttivo è il compito di selezione di Wason (1966). \\
	{\color{red}Approfondimento su slide "Intro\_IA-02-Generalità" pagina 99}
	%Argomento approfondito \hyperlink{90}{più avanti}
	\subsubsection{Approccio sub-simbolico}
	I modelli di IA in questo paradigma sono rappresentati implicitamente. \\
	La rappresentazione implicita deriva dall'apprendimento dall'esperienza senza rappresentazione simbolica di regole e proprietà. \\ 
	Invece di relazioni leggibili dall'uomo chiaramente definite, si implementano equazioni matematiche meno spiegabili per risolvere i problemi.\\
	Al fine di realizzare macchine che possano dirsi intelligenti, è necessario dotarle della capacità di estendere la propria conoscenza e le proprie abilità in modo autonomo.\\ \\
	\textbf{Campo machine learning}: Apprendimento automatico, costituito da diversi componenti basati su tecniche e paradigmi quali statistica, probabilità, reti, etc.\\
	Usato approccio bottom up (regole imparate da basso), in genere tratta di sistemi complessi.
	\paragraph{Ragionamento e apprendimento}
	La capacità di apprendimento è fondamentale in un sistema intelligente per:
	\begin{itemize}
		\item Risolvere nuovi problemi (e includere la scoperta nella propria base di conoscenza)
		\item Non ripetere gli errori fatti in passato
		\item Risolvere problemi in modo migliore o più efficiente
		\item Avere autonomia nell'affrontare e risolvere problemi
		\item Adattarsi ai cambiamenti dell'ambiente circostante
	\end{itemize}
	\textbf{Apprendimento induttivo}: Viene utilizzato per apprendere dei concetti espressi in genere con delle regole la cui forma è: "nella situazione X esegui l'azione Y", "all'osservazione X associa l'osservazione Y", "alla situazione X associa la conseguenza Y".\\
	Si parla di \textbf{generalizzazione} delle informazioni contenute nei pochi esempi a disposizione in un modello riutilizzabile in tutte altre possibili situazioni di quel tipo.
	\subsection{Cos'è la teoria del learning}
	\textbf{Sistemi esperti}: Un sistema esperto è un sistema basato su regole di produzione (della forma IF x THEN y, syllogismo) che, a partire da alcuni fatti (dati) cerca di dimostrare un'ipotesi o raggiungere una conclusione.\\ \\
	\textbf{Alberi di decisione}: I nodi degli alberi costituiscono le variabili, le foglie le decisioni.
	\begin{figure}[h]
		\centering
		\adjustbox{cfbox=white 5pt 0cm}{\includegraphics[scale=0.3]{example-image-golden}}
		\caption{Albero di decisione}
		\label{fig:im-3}
	\end{figure}
	\begin{center}
		\begin{tabular}{ |c|c|c|c| }
			\hline
			\textbf{Weather} & \textbf{Wind} & \textbf{Humidity} & \textbf{Decision} \\
			\hline
			Sunny & NO & 30-60 & PLAY\\
			\hline
			Sunny & YES & 90 & DON'T PLAY\\
			\hline
			Overcast & NO & 50 & PLAY\\
			\hline
			Rainy & YES & 50 & DON'T PLAY\\
			\hline
			Rainy & YES & 90 & DON'T PLAY\\
			\hline
			Sunny & NO & 80 & DON'T PLAY\\
			\hline
			Overcast & YES & 80 & PLAY\\
			\hline
			Rainy & NO & 70 & PLAY\\
			\hline
		\end{tabular}
	\end{center}
	Si parla di "Machine Learning simbolico" cioè basato su concetti astratti e non rappresentati da numeri o misure.\\
	Per creare gli alberi di decisione, solitamente si costruisce una “funzione/tabella” di verità.
	\subsubsection{Inferenze}
	\textbf{Learning = Inferencing + Memorizing}\\
	Nel learning se il risultato di un'inferenza risulta utile, essa viene memorizzata.\\
	Le inferenze possono essere \textbf{deduttive} o \textbf{induttive}.\\
	Si consideri l'equazione fondamentale per l'inferenza:
	\begin{itemize}
		\item $P$ sta per premessa;
		\item $BK$ per conoscenza pregressa;
		\item $\rightmodels$ per implicazione; %Dal pachetto MnSymbol
		\item $C$ per conseguente.
	\end{itemize}
	L'inferenza deduttiva deriva $C$, dai dati $P$ e $BK$ (veri), e preserva la verità. \\
	L'inferenza induttiva sta ipotizzando $P$, dai dati $C$ e $BK$, e preserva la falsità.\\ \\
	Si possono poi classificare ulteriormente in inferenza conclusiva (forte) e contingente (debole).
	\begin{itemize}
		\item \textbf{Inferenze conclusive}: Implicano regole di inferenza indipendenti dal dominio (contesto), mentre le inferenze contingenti implicano regole dipendenti dal dominio.
		\item \textbf{Deduzione contingente}: Produce probabili conseguenze di determinate cause.
		\item \textbf{Induzione contingente}: Produce probabili cause di determinate conseguenze.
		\item \textbf{Analogia}: area centrale del diagramma che rappresenta i processi di learning caratterizzati da induzione e deduzione.
	\end{itemize}
	\begin{figure}[h]
		\centering
		\adjustbox{cfbox=white 5pt 0cm}{\includegraphics[scale=0.3]{example-image-golden}}
		\caption{Schema riassuntivo}
		\label{fig:im-4}
	\end{figure}
	\subsubsection{Deep Learning e Reti Neurali Artificiali (ANN)}
	Il processo di deep learning è uno degli approcci all'apprendimento automatico che ha preso spunto dalla struttura del cervello, ovvero l'interconnessione dei vari neuroni.\\ \\
	Le reti neurali artificiali (ANN) si basano su questo fenomeno biologico, ma sono formate da neuroni artificiali costituiti da moduli software chiamati nodi. \\
	I nodi utilizzano calcoli matematici (al posto dei segnali chimici cerebrali) per comunicare e trasmettere informazioni. Questa rete neurale simulata elabora i dati raggruppando punti dati e facendo previsioni. \\ \\
	\textbf{Deep learning}: Si può immaginare come un diagramma di flusso con un livello iniziale in cui vengono inseriti dati e uno finale che fornisce risultati. \\
	Tra questi due troviamo "livelli nascosti" che elaborano informazioni a diversi livelli, modificando e adeguando il loro comportamento man mano che ricevono nuovi dati.\\ \\
	Alcuni limiti:
	\begin{itemize}
		\item Hanno bisogno di un numero enorme di esempi 
		\item Hanno difficoltà a generalizzare
		\item Hanno scarse capacità a gestire l’inferenza e il ragionamento logico
		\item Hanno difficoltà a distinguere fra correlazioni e causa-effetto.
	\end{itemize}
	\subsection{Black Box AI \& Explainable AI}
	I sistemi Machine Learning spesso funzionano come \textbf{una scatola nera}, che fa previsioni e decisioni come fanno gli umani, ma senza essere in grado di comunicare le ragioni per farlo.\\ \\
	\textbf{Explainable AI (XAI)}: Studia il motivo per cui è stata presa una decisione dall’AI in modo che i modelli di AI possano essere più interpretabili per gli utenti umani e consentire loro di capire perché il sistema è arrivato a una decisione specifica.
	\subsubsection{Indagare un sistema di AI}
	Un primo livello di indagine riguarda l’\textbf{Interpretability} (interpretabilità), cioè la possibilità di mettere in relazione causale i dati in ingresso con quelli in uscita (corrisponde alla previsione WHAT).\\
	Un secondo livello è relativo alla \textbf{Explainability} (spiegazione), in termini comprensibili ad un essere umano, su come il modello sia arrivato ad una determinata scelta o previsione (WHY).\\ \\
	La interpretability di un modello può generare la rappresentazione di un processo decisionale, ma trasformarlo in una explainability utile può essere difficile poiché dipende da diversi fattori:
	\begin{itemize}
		\item Il tipo di algoritmo che genera il modello
		\item Il livello di spiegazione richiesto
		\item Il tipo di dati utilizzati nel modello.
	\end{itemize}
	Alcuni approcci per la giustificazione a posteriori di un modello di ML:
	\begin{itemize}
		\item Spiegazione testuale
		\item Spiegazione tramite esempi
		\item Semplificazione del modello
		\item Visualizzazione
		\item Spiegazione locale: Riduzione della complessità del problema grazie alla ricerca di spiegazioni per una sotto-porzione del dominio approssimato dal modello.
		\item Feature rilevanti: Quantificazione dell'influenza di una feature del problema sull'output finale. L'analisi comparata di questi attributi fornisce una spiegazione indiretta sul funzionamento del modello.
	\end{itemize}
	\begin{figure}[h]
		\centering
		\adjustbox{cfbox=white 5pt 0cm}{\includegraphics[scale=0.3]{example-image-golden}}
		\caption{Schema riassuntivo}
		\label{fig:im-5}
	\end{figure}
	\paragraph{Esempio: Un oggetto cade a terra e si rompe.}
	\begin{itemize}
		\item L’Interpretability può dare le ragioni meccaniche della caduta ma senza coglierne veramente le cause: l’oggetto è caduto, per questo ha toccato terra e per tale ragione si è infranto.
		\item L’Explainability dovrebbe dare le vere ragioni alla base del fenomeno. In questo caso potrebbe essere: l’oggetto è caduto, a causa della gravità ha ottenuto una certa accelerazione e l’energia accumulata nella caduta ha portato alla rottura dell’oggetto all’impatto.
	\end{itemize}
\end{document}